% Шаблон отчёта. Компилировать: xelatex report.tex (дважды, ради содержания)
\documentclass[a4paper,14pt]{extarticle}
\usepackage{../../../template/gostreport}

\begin{document}

\makegosttitle
  {1}                                   % номер работы
  {Кодирование и шифрование}            % название работы
  {студент группы R3243 Осипов А.\,П.\\[4pt]
   студент группы R???? Глазков В.\,?.\\[4pt]
   студент группы R???? Прохоревич И.\,?.}  % исполнители, по строке на человека
  {Догадин Е.\,В.}                      % преподаватель

\tableofcontents
\newpage

\unnumsection{Введение}

Кратко: что за работа, какие инструменты применяются, какая цель.

\section{Задание 1. Название}

Текст, расчёты, обоснования.

\begin{equation}
  \mat{A} = \begin{pmatrix} a & b \\ c & d \end{pmatrix},
  \qquad \det \mat{A} = ad - bc.
\end{equation}

Результат преобразования показан на рисунке~\ref{fig:example}.

\begin{figure}[H]
  \centering
  \includegraphics[width=0.7\textwidth]{../figures/example.pdf}
  \caption{Пример подписи к рисунку}
  \label{fig:example}
\end{figure}

\subsection*{Вывод по заданию 1}

Не пересказ действий, а анализ: почему результат такой, что он означает.

\unnumsection{Заключение}

Общий вывод по работе.

\newpage
\unnumsection{Приложение А. Код}

\lstinputlisting[style=py,caption={Решение задания 1},label={lst:task1}]{../code/task1.py}

\end{document}
